% SPDX-License-Identifier: CC0-1.0
%
% Pseudocode of the procedures that decideHamPath runs, transcribed from the Lean definitions.
% Requires in the preamble of hpst.tex:
%   \usepackage{algorithm}
%   \usepackage[noend]{algpseudocode}
% Uses \lean{...} from hpst.tex.
%
% Procedures are typeset in small capitals (\proc), helper predicates and functions in sans
% serif (\pfn), multi-letter variables in italics (\var).
%
\algrenewcommand{\textproc}{\textsc}
\algrenewcommand{\algorithmiccomment}[1]{\hfill{\small$\triangleright$~#1}}
\newcommand{\proc}[1]{\mbox{\textsc{#1}}}
\newcommand{\pfn}[1]{\textsf{#1}}
\newcommand{\var}[1]{\mathit{#1}}
\newcommand{\pctrue}{\textsf{true}}
\newcommand{\pcfalse}{\textsf{false}}
\newcommand{\pcnone}{\textsf{none}}
\newcommand{\pcbreak}{\textbf{break}}

\paragraph{Pseudocode.}
\proc{DecideHamPath} (\cref{alg:decideHamPath}) chooses one of three procedures by testing which
endpoints are inner.
Each walk deletes endpoints until the two out-vertices are Far with phantoms that are not opposed,
and then runs a final test.
The final test builds a path factor with \proc{FindFactor}, in the hole cases moves it into each
class with \proc{ClassSearch}, and runs the two-phantom loop \proc{TwoPhantomLoop}.
Each round of that loop is a \proc{RunPhase} followed by \proc{ReducePieces}.

The blocks below use the following conventions.
Each procedure has the name of its Lean definition, capitalised, and the caption of its block gives
that name.
All of these definitions are in \lean{Algorithm.lean}.
Helper predicates and functions are set in sans serif and are described after each block with
their Lean names.
$\pfn{Nbrs}(t)$ (Lean \lean{nbrsL}) is the list of pairs $(y,e)$ with $y$ a lattice neighbour of
$t$ and $e$ the edge $t$--$y$, in the order right, left, up, down.
The cells of $V$ are visited in lexicographic order (\lean{cellList}).
The procedures return Booleans, and \pctrue\ means that a Hamiltonian path exists.
Every loop over a list stops at its first success, as \lean{List.any}, \lean{List.find?} and
\lean{List.findSome?} do in Lean.
A loop ``for $n$ rounds'' runs at most $n$ times, and $n$ is the fuel of the Lean definition.

\begin{algorithm}[htbp]
\caption{\proc{DecideHamPath} (Lean \lean{decideHamPath})}
\label{alg:decideHamPath}
\begin{algorithmic}[1]
\Function{DecideHamPath}{$V,s,t$}
  \If{$\pfn{IsInner}(V,s)$}
    \If{$\pfn{IsInner}(V,t)$} \Return \Call{TwoHoleCase}{$V,s,t$}
      \Comment{both endpoints inner}
    \Else\ \Return \Call{OneHoleCase}{$V,t,s$}
      \Comment{only $s$ inner, read the path backwards}
    \EndIf
  \ElsIf{$\pfn{IsInner}(V,t)$} \Return \Call{OneHoleCase}{$V,s,t$}
    \Comment{only $t$ inner}
  \Else\ \Return \Call{BoundaryEndpoints}{$V,s,t$}
    \Comment{neither endpoint inner}
  \EndIf
\EndFunction
\end{algorithmic}
\end{algorithm}

In \cref{alg:decideHamPath}, $\pfn{IsInner}(V,p)$ (Lean \lean{isInner}) holds when the four squares with
corner $p$ are cells of $V$.

\begin{algorithm}[htbp]
\caption{\proc{BoundaryEndpoints} (Lean \lean{boundaryEndpoints})}
\label{alg:boundaryEndpoints}
\begin{algorithmic}[1]
\Function{BoundaryEndpoints}{$V,s,t$}
  \Comment{$s$ and $t$ each have a missing king neighbour}
  \If{$|V|\le2$} \Return [$s$ is a lattice neighbour of $t$]
  \EndIf
  \If{some $(q,e_{sq})\in\pfn{Nbrs}(s)$ has $q\notin V$ (take the first)}
    \Comment{$q$ becomes the phantom of $s$}
    \State \Return \Call{BoundaryEndpointsWalk}{$s,q,e_{sq},12,V,t$}
  \EndIf
  \ForAll{$(x,e)\in\pfn{Nbrs}(s)$ with $x\in V$, $x\ne t$}
    \Comment{delete $s$ (phantom of $x$)}
    \If{\Call{BoundaryEndpointsWalk}{$x,s,e,12,V\setminus\{s\},t$}} \Return \pctrue
    \EndIf
  \EndFor
  \State \Return \pcfalse
\EndFunction
\end{algorithmic}
\end{algorithm}

\begin{algorithm}[htbp]
\caption{\proc{BoundaryEndpointsWalk} (Lean \lean{boundaryEndpointsWalk})}
\label{alg:boundaryEndpointsWalk}
\begin{algorithmic}[1]
\Function{BoundaryEndpointsWalk}{$s,q,e_{sq},\var{fuel},V,t$}
  \Comment{$s$ keeps its phantom $q$, $t$ moves}
  \If{$\var{fuel}=0$} \Return \pcfalse
  \EndIf
  \If{$|V|\le2$} \Return [$s$ is a lattice neighbour of $t$]
  \EndIf
  \If{some $(p,e_{tp})\in\pfn{Nbrs}(t)$ has $\pfn{IsUsablePhantom}(V,s,q,t,p)$ (take the first)}
    \If{$\pfn{IsFar}(t,s)$}
      \Comment{$p$ is usable and $t$ is Far from $s$}
      \State \Return \Call{TwoPhantomTest}{$V,s,t,p,q,e_{tp},e_{sq}$}
      \Comment{the final test}
    \EndIf
  \EndIf
  \ForAll{$(x,\cdot)\in\pfn{Nbrs}(t)$ with $x\in V$, $x\ne s$}
    \Comment{delete $t$, $x$ is the next endpoint}
    \If{\Call{BoundaryEndpointsWalk}{$s,q,e_{sq},\var{fuel}-1,V\setminus\{t\},x$}} \Return \pctrue
    \EndIf
  \EndFor
  \State \Return \pcfalse
\EndFunction
\end{algorithmic}
\end{algorithm}

In \cref{alg:boundaryEndpoints,alg:boundaryEndpointsWalk}, $e_{sq}$ is the phantom edge $s$--$q$ and $e_{tp}$ is the phantom
edge $t$--$p$.
The fuel $12$ allows the at most $11$ deletion levels of \cref{sec:walk}.
$\pfn{IsUsablePhantom}(V,s,q,t,p)$ (Lean \lean{isUsablePhantom}) holds when $p\notin V$, $p\ne q$, and
the phantom $p$ of $t$ is not opposed to the phantom $q$ of $s$, so $p$ is a usable phantom.
$\pfn{IsFar}(t,s)$ (Lean \lean{isFar}) holds when $t=s$ or $\dist(t,s)\ge2$.

\begin{algorithm}[htbp]
\caption{\proc{HoleCase} (Lean \lean{holeCase}), \proc{OneHoleCase} (Lean
  \lean{oneHoleCase}) and \proc{TwoHoleCase} (Lean \lean{twoHoleCase})}
\label{alg:holeCase}
\label{alg:oneHoleCase}
\label{alg:twoHoleCase}
\begin{algorithmic}[1]
\Function{HoleCase}{$\var{finalTest},V,s,t$}
  \Comment{$t$ is inner}
  \If{$|V|\le2$} \Return [$s$ is a lattice neighbour of $t$]
  \EndIf
  \ForAll{$(w,e_{wt})\in\pfn{Nbrs}(t)$ with $w\in V$, $w\ne s$}
    \Comment{delete $t$ (phantom of $w$)}
    \If{$|V\setminus\{t\}|\le2$}
      \If{$s$ is a lattice neighbour of $w$} \Return \pctrue
      \EndIf
    \Else
      \ForAll{$(x,e)\in\pfn{Nbrs}(s)$ with $x\in V\setminus\{t\}$, $x\ne w$}
        \Comment{delete $s$ (phantom of $x$)}
        \If{\Call{HoleCaseWalk}{$\var{finalTest}(w,e_{wt}),w,t,10,V\setminus\{t,s\},x,s,e$}}
          \Return \pctrue
        \EndIf
      \EndFor
    \EndIf
  \EndFor
  \State \Return \pcfalse
\EndFunction
\Statex
\Function{OneHoleCase}{$V,s,t$}
  \Comment{$t$ inner, $s$ not inner}
  \State $\var{finalTest}(w,e_{wt})(U,s',q,e_{sq})\gets\Call{OneHoleFinalTest}{U,s',w,t,q,e_{wt},e_{sq}}$
  \State \Return \Call{HoleCase}{$\var{finalTest},V,s,t$}
\EndFunction
\Statex
\Function{TwoHoleCase}{$V,s,t$}
  \Comment{the original $s$ is the second hole}
  \State $\var{finalTest}(w,e_{wt})(U,s',q,e_{sq})\gets\Call{TwoHoleFinalTest}{U,s',w,t,s,q,e_{wt},e_{sq}}$
  \State \Return \Call{HoleCase}{$\var{finalTest},V,s,t$}
\EndFunction
\end{algorithmic}
\end{algorithm}

In \cref{alg:holeCase} the point $w$ is the neighbour of $t$ on the path, $e_{wt}$ is the edge
$w$--$t$, and the deleted $t$ is the phantom of $w$.
The argument $\var{finalTest}$ is a function of $(w,e_{wt})$ that returns the final test of the
walk.
In \proc{TwoHoleCase} the original $s$ is passed to \proc{TwoHoleFinalTest} as the second hole
$r$.

\begin{algorithm}[htbp]
\caption{\proc{HoleCaseWalk} (Lean \lean{holeCaseWalk})}
\label{alg:holeCaseWalk}
\begin{algorithmic}[1]
\Function{HoleCaseWalk}{$\var{finalTest},w,p_w,\var{fuel},U,s,q,e_{sq}$}
  \Comment{$w$ keeps its phantom $p_w$, $s$ moves}
  \If{$\var{fuel}=0$} \Return \pcfalse
  \EndIf
  \If{$|U|\le2$} \Return [$s$ is a lattice neighbour of $w$]
  \EndIf
  \If{$\pfn{IsFar}(w,s)$ and not $\pfn{PhantomsOpposed}(w,p_w,s,q)$}
    \Comment{ready for the final test}
    \State \Return $\var{finalTest}(U,s,q,e_{sq})$
  \EndIf
  \ForAll{$(x,e)\in\pfn{Nbrs}(s)$ with $x\in U$, $x\ne w$}
    \Comment{delete $s$ (phantom of $x$)}
    \If{\Call{HoleCaseWalk}{$\var{finalTest},w,p_w,\var{fuel}-1,U\setminus\{s\},x,s,e$}}
      \Return \pctrue
    \EndIf
  \EndFor
  \State \Return \pcfalse
\EndFunction
\end{algorithmic}
\end{algorithm}

In \cref{alg:holeCaseWalk} the phantom $p_w$ of $w$ is the deleted inner point $t$.
The phantom $q$ of $s$ is the point deleted last, and $e_{sq}$ is the edge $s$--$q$.
The fuel $10$ allows the at most $9$ deletions of \cref{sec:one}.
$\pfn{PhantomsOpposed}(a,m_a,b,m_b)$ (Lean \lean{phantomsOpposed}) holds when the phantoms $m_a$ of $a$ and
$m_b$ of $b$ are opposed (\cref{def:opp}).

\begin{algorithm}[htbp]
\caption{\proc{OneHoleFinalTest} (Lean \lean{oneHoleFinalTest}) and
  \proc{TwoHoleFinalTest} (Lean \lean{twoHoleFinalTest})}
\label{alg:oneHoleFinalTest}
\label{alg:twoHoleFinalTest}
\begin{algorithmic}[1]
\Function{OneHoleFinalTest}{$U,s,w,t,q,e_{wt},e_{sq}$}
  \Comment{hole at $t$}
  \If{not $\pfn{IsNarrow}(U)$} \Return \pcfalse
    \Comment{a wide set has no Hamiltonian path}
  \EndIf
  \State $R\gets\Call{FindFactor}{U,d_{sw}}$
  \Comment{a path factor with ends $s$ and $w$}
  \If{$R=\pcnone$} \Return \pcfalse
  \EndIf
  \State $F\gets R\cup\{e_{wt},e_{sq}\}$
  \Comment{the reference factor, phantoms $t$ and $q$}
  \For{$h=-|U|,\dots,|U|$}
    \Comment{every possible class (\lean{classRange})}
    \State $G\gets\Call{ClassSearch}{U,F,[t],[h]}$
    \Comment{a factor of class $h$ (\lean{classSearchOne})}
    \If{$G\ne\pcnone$ and \Call{TwoPhantomLoop}{$U,t,q,|U|+2,G$}} \Return \pctrue
    \EndIf
  \EndFor
  \State \Return \pcfalse
\EndFunction
\Statex
\Function{TwoHoleFinalTest}{$U,s,w,t,r,q,e_{wt},e_{sq}$}
  \Comment{holes at $t$ and $r$}
  \If{not $\pfn{IsNarrow}(U)$} \Return \pcfalse
  \EndIf
  \State $R\gets\Call{FindFactor}{U,d_{sw}}$
  \If{$R=\pcnone$} \Return \pcfalse
  \EndIf
  \State $F\gets R\cup\{e_{wt},e_{sq}\}$
  \For{$h_1=-|U|,\dots,|U|$}
    \For{$h_2=-|U|,\dots,|U|$}
      \Comment{every possible pair of classes}
      \State $G\gets\Call{ClassSearch}{U,F,[t,r],[h_1,h_2]}$
      \If{$G\ne\pcnone$ and \Call{TwoPhantomLoop}{$U,t,q,|U|+2,G$}} \Return \pctrue
      \EndIf
    \EndFor
  \EndFor
  \State \Return \pcfalse
\EndFunction
\end{algorithmic}
\end{algorithm}

In \cref{alg:oneHoleFinalTest}, $\pfn{IsNarrow}(U)$ (Lean \lean{isNarrow}) holds when the width and the
height of the bounding box of $U$ are both less than $|U|$.
The degrees $d_{sw}$ are $1$ at $s$ and $w$ and $2$ elsewhere.
The factor $F$ is a two-phantom factor on $U\cup\{t,q\}$ with phantoms $t$ of $w$ and $q$ of $s$.

\begin{algorithm}[htbp]
\caption{\proc{ClassSearch} (Lean \lean{classSearch}, with \lean{classPotential} and
  \lean{relaxA})}
\label{alg:classSearch}
\begin{algorithmic}[1]
\Function{ClassSearch}{$V,F,[a_1,\dots,a_k],[h_1,\dots,h_k]$}
  \Comment{hole $a_i$ gets class $h_i$}
  \State $\var{box}\gets\pfn{BoxSquares}(V)$
  \Comment{squares of the bounding box, $x$-major}
  \ForAll{$c\in\var{box}$}
    \Comment{start below every solution (\lean{kInitA})}
    \State $\varphi(c)\gets h_i$ for the first $i$ whose hole cell $a_i$ is $c$, else $-\pfn{BoxWidth}(V)$
  \EndFor
  \State $\varphi(c)\gets0$ for every square $c\notin\var{box}$
  \For{$|\var{box}|+1$ rounds}
    \Comment{Bellman--Ford relaxation (\lean{relaxA})}
    \State $\varphi'(c)\gets\max(\varphi(c),\pfn{LowerBound}(V,F,\varphi,c))$ for all
      $c\in\var{box}$
    \If{$\varphi'=\varphi$ on $\var{box}$} \pcbreak
      \Comment{fixed point, the least potential}
    \EndIf
    \State $\varphi\gets\varphi'$
  \EndFor
  \If{$\pfn{IsAdmissible}(V,F,[a_i],[h_i],\varphi)$}
    \Return $F\oplus\pfn{JumpEdges}(V,\varphi)$
    \Comment{the factor of these classes}
  \Else\ \Return \pcnone
    \Comment{no factor has these classes}
  \EndIf
\EndFunction
\end{algorithmic}
\end{algorithm}

In \cref{alg:classSearch}, $\pfn{BoxSquares}(V)$ (Lean \lean{boxSquares}) lists the squares of the bounding
box of $V$, and $\pfn{BoxWidth}(V)$ (Lean \lean{boxWidth}) is the width of that box.
Let $\sigma(e)$ be the sign of the arc of $F$ across $e$ ($+1$ from the lower or left square to the
upper or right one) when $e$ is an edge of $V$, and $0$ otherwise.
$\pfn{LowerBound}(V,F,\varphi,c)$ (Lean \lean{lowerBound}) is the largest lower bound on $\varphi(c)$
given by the four neighbours of $c$ under the constraint that each jump lies between $0$ and
$\sigma(e)$.
It is the maximum of $\varphi(c_\downarrow)+\min(0,\sigma)$ and
$\varphi(c_\leftarrow)+\min(0,\sigma)$ across the bottom and left sides and of
$\varphi(c_\uparrow)-\max(0,\sigma)$ and $\varphi(c_\rightarrow)-\max(0,\sigma)$ across the top
and right sides.
$\pfn{IsAdmissible}$ (Lean \lean{isAdmissible}) holds when every side of every square of the box has jump
$0$, or is an edge $e$ of $V$ with jump $\sigma(e)$, and $\varphi(a_i)=h_i$ for all $i$.
$\pfn{JumpEdges}(V,\varphi)$ (Lean \lean{jumpEdges}) is the set of edges of $V$ on sides of squares of the
box across which $\varphi$ jumps.
The rounds bound $|\var{box}|+1$ is \lean{kFuel}.
The potentials are stored in an array indexed by box position (\lean{toFA}).
With one hole, \proc{ClassSearch} is the Lean \lean{classSearchOne}.

\begin{algorithm}[htbp]
\caption{\proc{TwoPhantomTest} (Lean \lean{twoPhantomTest}) and
  \proc{TwoPhantomLoop} (Lean \lean{twoPhantomLoop})}
\label{alg:twoPhantomTest}
\label{alg:twoPhantomLoop}
\begin{algorithmic}[1]
\Function{TwoPhantomTest}{$V,s,t,p,q,e_{tp},e_{sq}$}
  \Comment{phantoms $p$ of $t$ and $q$ of $s$}
  \State $R\gets\Call{FindFactor}{V,d_{st}}$
  \Comment{a path factor with ends $s$ and $t$}
  \If{$R=\pcnone$} \Return \pcfalse
  \EndIf
  \State \Return \Call{TwoPhantomLoop}{$V,p,q,|V|+2,R\cup\{e_{tp},e_{sq}\}$}
  \Comment{add the phantom edges}
\EndFunction
\Statex
\Function{TwoPhantomLoop}{$V,p,q,\var{fuel},G$}
  \State $W\gets V\cup\{p,q\}$
  \For{$\var{fuel}$ rounds}
    \Comment{each round joins two pieces}
    \If{$\ncomp(W,G)=1$} \Return \pctrue
      \Comment{one piece, a Hamiltonian path}
    \EndIf
    \State $(H,\var{bridges},\var{outcome})\gets\Call{RunPhase}{V,W,|W|+1,G,[\,]}$
    \If{$\var{outcome}\in\{\textsf{stuck},\textsf{fuel}\}$} \Return \pcfalse
    \EndIf
    \State $G\gets\Call{ReducePieces}{V,W,H,\var{bridges},\var{outcome}}$
    \Comment{one piece fewer, no bridges}
  \EndFor
  \State \Return $[\ncomp(W,G)=1]$
\EndFunction
\end{algorithmic}
\end{algorithm}

\begin{algorithm}[htbp]
\caption{\proc{RunPhase} (Lean \lean{runPhase})}
\label{alg:runPhase}
\begin{algorithmic}[1]
\Function{RunPhase}{$V,W,\var{fuel},H,\var{bridges}$}
  \Comment{$\var{bridges}$ is listed newest first}
  \For{$\var{fuel}$ steps}
    \State $\var{conn}\gets\pfn{ConnectedPairs}(W,H)$
    \Comment{all pairs of points in one piece}
    \State $c\gets$ the first merging cell of $V$ that is not a pair square
    \Comment{\lean{findMergingCell}}
    \If{$c\ne\pcnone$} \Return $(H,\var{bridges},\textsf{merge}\ c)$
      \Comment{flipping $c$ joins two pieces}
    \EndIf
    \State $b\gets\pfn{FindRow}(V,\var{conn},H,\var{bridges})$
    \Comment{a critical cell whose row is not stuck}
    \If{$b=\pcnone$} \Return $(H,\var{bridges},\textsf{stuck})$
    \EndIf
    \If{$r_b$ is odd} \Return $(H,\var{bridges},\textsf{odd}\ b)$
      \Comment{lowering the odd row joins two pieces}
    \EndIf
    \If{$\pfn{FlipSplits}(W,H,b)$}
      \Comment{bridge step, undone or lowered later}
      \State flip $E_b$ in $H$, \quad $\var{bridges}\gets b::\var{bridges}$
      \Comment{one piece more}
    \Else\ \Return $(H,\var{bridges},\textsf{twist}\ b)$
      \Comment{even row, flipping $E_b$ splits nothing}
    \EndIf
  \EndFor
  \State \Return $(H,\var{bridges},\textsf{fuel})$
\EndFunction
\end{algorithmic}
\end{algorithm}

In \cref{alg:runPhase}, a bridge $b=(q_b,a_b,r_b)$ is a row of length $r_b$ from the cell $q_b$ in
the frame with front $a_b$, and $E_b$ is its last cell $c_{r_b-1}$.
Its fixed pair is $(c_{r_b-2},c_{r_b-1})$, and these two cells are \emph{pair squares} (Lean
\lean{pairSqB}).
$\pfn{ConnectedPairs}(W,H)$ (Lean \lean{connectedPairs}) is the set of pairs of points of $W$ that lie in one
piece of $H$.
A merging cell is an alternating cut cell (\lean{Alt}, \lean{isCutCell}).
$\pfn{FindRow}$ (Lean \lean{findRow}) scans the cells $q$ of $V$ in order and for each the sides
$a$ in the order bottom, top, left, right.
It returns $(q,a,r)$ for the first $q$ that is critical in frame $a$ with dark back and light front
(\lean{isCritical}), is not a square of any bridge, and has a row of length $r$ that is not stuck
(\lean{scanRow}).
\lean{scanRow} steps $j=1,2,\dots$ up to $|V|+1$ times.
It ends with a row of length $j$ when the front arc of $c_{j-1}$ differs from that of $q$, and it
ends stuck when $c_j$ is not a cell of $V$ or is a cut cell.
$\pfn{FlipSplits}(W,H,b)$ (Lean \lean{flipSplits}) holds when, after the flip of $E_b$, the corners
$b_1$ and $f_1$ of $E_b$ lie in different pieces.

\begin{algorithm}[htbp]
\caption{\proc{ReducePieces} (Lean \lean{reducePieces}) and
  \proc{UndoBridges} (Lean \lean{applyAll} and \lean{undoPlan})}
\label{alg:reducePieces}
\begin{algorithmic}[1]
\Function{ReducePieces}{$V,W,H,[b_1,\dots,b_k],\var{outcome}$}
  \Comment{$b_1$ is the newest bridge}
  \State $H_0\gets H$
  \For{$i=1,\dots,k$}
    \Comment{the states before each bridge (\lean{bridgeStatesConn})}
    \State $H_i\gets H_{i-1}$ with $E_{b_i}$ flipped, \quad $\var{conn}_i\gets\pfn{ConnectedPairs}(W,H_i)$
  \EndFor
  \If{$\var{outcome}=\textsf{merge}\ c$}
    \State $X\gets$ the first cell of $V$ with
      $\pfn{IsGoodMergingCell}(W,H,[b_i],[\var{conn}_i],X)$, else $c$
    \Comment{Theorem~R}
    \State $(u,v)\gets\pfn{MergePair}(H,X)$, \quad flip $X$ in $H$
    \Comment{half bridge}
  \ElsIf{$\var{outcome}=\textsf{odd}\ b$ or $\var{outcome}=\textsf{twist}\ b$}
    \State $(u,v)\gets\pfn{RowStartPair}(b)$
    \State flip the cells $c_{r_b-1},\dots,c_1,c_0$ of the row of $b$ in $H$
    \Comment{half bridge, lower the row}
  \Else\ \Return $H$
  \EndIf
  \State \Return \Call{UndoBridges}{$H,[b_1,\dots,b_k],[\var{conn}_1,\dots,\var{conn}_k],(u,v)$}
\EndFunction
\Statex
\Function{UndoBridges}{$H,[b_1,\dots,b_k],[\var{conn}_1,\dots,\var{conn}_k],(u,v)$}
  \Comment{$(u,v)$ is the pending pair}
  \For{$i=1,\dots,k$}
    \Comment{the plan (\lean{undoPlan})}
    \If{$u$ and $v$ are connected in $\var{conn}_i$}
      \Comment{joined before $b_i$ was made}
      \State $\var{op}_i\gets$ lower $b_i$, \quad $(u,v)\gets\pfn{RowStartPair}(b_i)$
    \Else\ $\var{op}_i\gets$ undo $b_i$
    \EndIf
  \EndFor
  \For{$i=1,\dots,k$}
    \Comment{apply the plan (\lean{applyAll})}
    \If{$\var{op}_i$ lowers $b_i$} flip the cells $c_{r_{b_i}-2},\dots,c_0$ of $b_i$ in $H$
    \Else\ flip $E_{b_i}$ in $H$
      \Comment{back to the state before $b_i$}
    \EndIf
  \EndFor
  \State \Return $H$
\EndFunction
\end{algorithmic}
\end{algorithm}

In \cref{alg:reducePieces}, $\var{conn}_i$ is the connection set of the state before bridge $b_i$
was made.
The half bridge joins the pieces of $u$ and $v$, and \proc{UndoBridges} removes the bridges while
keeping that join.
$\pfn{MergePair}(H,X)$ (Lean \lean{mergePair}) is the pair $(b_1,f_1)$ of $X$ in the frame with front
top when the bottom of $X$ is dark and front right otherwise (\lean{altFrame}).
$\pfn{RowStartPair}(b)$ (Lean \lean{rowStartPair}) is the pair $(b_1,f_1)$ of $q_b$ in the frame
$a_b$.
$\pfn{IsGoodMergingCell}$ (Lean \lean{isGoodMergingCell}) tests that $X$ is alternating, cut, and not a pair
square, and that every bridge $b_i$ whose square beyond $E_{b_i}$ is $X$ is lowered by the plan
started from $\pfn{MergePair}(H,X)$, so $X$ is good in the sense of Theorem~R.
The plan and the operations are computed and applied in the order $b_1,\dots,b_k$, newest bridge
first.

\begin{algorithm}[htbp]
\caption{\proc{FindFactor} (Lean \lean{findFactor} with \lean{augLoop}, \lean{augStep} and
  \lean{layerAcc}) and \proc{AugmentingWalk} (Lean \lean{walk})}
\label{alg:findFactor}
\begin{algorithmic}[1]
\Function{FindFactor}{$V,d$}
  \Comment{a factor with degrees $d$, or none}
  \State $F\gets\emptyset$
  \For{$|V|+1$ rounds}
    \Comment{one augmentation per round (\lean{augLoop})}
    \State $\var{layer}_0\gets\{z\in V : z \text{ black},\ \deg_F z<d(z)\}$
    \Comment{deficient black points}
    \For{$k=1,\dots,|V|$}
      \Comment{\lean{layerAcc}}
      \State $\var{layer}_k\gets\var{layer}_{k-1}\cup\pfn{AltNbrs}(V,F,\var{layer}_{k-1})$
    \EndFor
    \State $y\gets$ the least white $y\in\var{layer}_{|V|}$ with $\deg_F y<d(y)$
    \Comment{a deficient white end}
    \If{no such $y$} \pcbreak
      \Comment{no augmenting walk is left}
    \EndIf
    \State $F\gets\Call{AugmentingWalk}{V,F,[\var{layer}_{|V|-1},\dots,\var{layer}_0],y}$
  \EndFor
  \If{$\deg_F z=d(z)$ for all $z\in V$} \Return $F$
  \Else\ \Return \pcnone
  \EndIf
\EndFunction
\Statex
\Function{AugmentingWalk}{$V,F,[L_{n-1},\dots,L_0],v$}
  \Comment{walk from $v$ back to a black start}
  \State $F'\gets F$
  \For{$k=n-1,\dots,0$}
    \If{$v\notin L_k$}
      \Comment{step back into layer $k$}
      \State $e\gets$ the first edge $e$ at $v$ with $\pfn{IsBackEdge}(V,F,L_k,v,e)$
      \If{$e=\pcnone$} \Return $F'$
      \EndIf
      \State $F'\gets F'\setminus\{e\}$ if $v$ is black, else $F'\cup\{e\}$
      \Comment{toggle $e$}
      \State $v\gets$ the other end of $e$
    \EndIf
  \EndFor
  \State \Return $F'$
\EndFunction
\end{algorithmic}
\end{algorithm}

In \cref{alg:findFactor}, a point $(x,y)$ is black when $x+y$ is even and white otherwise.
$\pfn{AltNbrs}(V,F,u)$ (Lean \lean{nbr}) is the set of points joined to $u$ by an edge of $V$ that
is not in $F$ when $u$ is black and in $F$ when $u$ is white.
So $\var{layer}_k$ is the set of points reached from a deficient black point by an alternating
walk of at most $k$ steps.
$\pfn{IsBackEdge}(V,F,L,v,e)$ (Lean \lean{Bwd}) holds when $e$ is an edge of $V$ at $v$ whose other
end is in $L$, with $e\in F$ when $v$ is black and $e\notin F$ when $v$ is white.
It reads the factor $F$ of the start of the round, not $F'$.
The edges at $v$ are tried in the order right, left, up, down.
Each round toggles the edges of one augmenting walk from $y$ back to $\var{layer}_0$, and the
rounds stop when no deficient white point is reachable.
